% Source: https://artofproblemsolving.com/community/c3h2227851_perfectabundantdeficient_numbers
% Stable identity: aops-topic-2227851-q01
For a positive integer $n$, let $\sigma(n)$ be the sum of all positive divisors of $n$. Call $n$ \emph{perfect} if $\sigma(n)=2n$, \emph{abundant} if $\sigma(n)>2n$ and \emph{deficient} if $\sigma(n)<2n$. There is no known formula listing all numbers of each kind, but each kind has simple infinite families.
\begin{enumerate}
\item[(a)] Prove that if $2^p-1$ is prime, then $2^{p-1}\left(2^p-1\right)$ is perfect.
\item[(b)] Prove that every prime power $q^k$ ($q$ prime, $k\ge1$) is deficient.
\item[(c)] Prove that every multiple $kn$ ($k\ge2$) of a perfect number $n$ is abundant, and deduce that $12,18,24,30,\ldots$ are abundant.
\end{enumerate}
